\section{Powiązane prace}
\label{sec:related-work}

\TODO{rozwinąć po przeglądzie literatury; punkty wyjścia poniżej.}

\begin{itemize}
  \item Łamigłówki, w tym wieże Hanoi, jako test modeli z rozumowaniem; model podaje całą listę ruchów w jednej
    odpowiedzi, bez narzędzi~\cite{shojaee2025illusion}.
  \item Polemika: limit długości odpowiedzi jako część przyczyny załamania wyników~\cite{lawsen2025comment}.
  \item Te same łamigłówki z narzędziami (interpreter, notatnik)~\cite{song2025thinking}.
  \item Model uczony ze wzmocnieniem sortuje akcjami porównania i zamiany na częściowo ukrytym stanie; najbliższa
    praca do naszej~\cite{yue2025earl}.
  \item \TODO{benchmarki tool callingu i prace o rozpoznawaniu algorytmu z ciągu operacji.}
\end{itemize}
